% Part: computability
% Chapter: computability-theory
% Section: halting-problem

\documentclass[../../../include/open-logic-section]{subfiles}

\begin{document}

\olfileid{cmp}{thy}{hlt}
\olsection{थांबण्याची समस्या}

रचनेवरूनच, सार्वत्रिक आंशिक संगणनक्षम फलन~$\fn{Un}(e,x)$ तेव्हाच
परिभाषित असते जेव्हा~$e$ ने संकेतबद्ध केलेल्या फलनाच्या आदान~$x$
वरील संगणनातून मूल्य मिळते. असे होते की नाही हे आपण ठरवू शकतो का,
हा नैसर्गिक प्रश्न आहे. प्रत्यक्षात ते ठरवता येत नाही. ट्यूरिंग
यंत्रांच्या संगणन-प्रतिमानात याचा अर्थ, दिलेले ट्यूरिंग यंत्र दिलेल्या
आदानावर थांबते का हे संगणनाने ठरवता येत नाही. म्हणून पुढील प्रमेयाला
``थांबण्याच्या समस्येची अनिर्णेयता'' म्हणतात. खाली आपण दोन सिद्धता
देऊ. पहिली आपल्या आधीच्या चर्चेचा धागा पुढे नेते, तर दुसरी अधिक थेट आहे.

\begin{thm}
\ollabel{thm:halting-problem}
पुढीलप्रमाणे असू द्या:
\[
h(e, x)  =
\begin{cases}
1 & \text{जर\/ $\fn{Un}(e, x)$ परिभाषित असेल तर} \\
0 & \text{अन्यथा.}
\end{cases}
\]
मग $h$ संगणनक्षम नाही.
\end{thm}

\begin{proof}
$h$ संगणनक्षम आहे असे समजा. मग सार्वत्रिक संगणनक्षम फलन परिभाषित
करता येईल, हे दाखवू. पुढीलप्रमाणे व्याख्या करा:
\[
\fn{Un'}(e,x) =
\begin{cases}
\fn{Un}(e,x) & \text{जर $h(e,x) = 1$ असेल तर} \\
0 & \text{अन्यथा.}
\end{cases}
\]
आता $\fn{Un'}(e, x)$ हे सर्वत्र परिभाषित फलन आहे आणि $h$
संगणनक्षम असेल तर तेही संगणनक्षम आहे. उदाहरणार्थ, आदिम पुनरावर्तन
वापरून $g$ ची व्याख्या पुढीलप्रमाणे करता येते:
\begin{align*}
g(0, e, x) & \simeq 0 \\
g(y+1, e, x) & \simeq \fn{Un}(e,x);
\end{align*}
मग
\[
\fn{Un'}(e,x) \simeq g(h(e,x),e,x).
\]
$\fn{Un'}(e,x)$ हे $\fn{Un}(e,x)$ परिभाषित असेल तेथे त्याच्याशी
सहमत असल्याने, $\fn{Un'}$ हे सर्वत्र परिभाषित असलेल्या आंशिक
संगणनक्षम फलनांसाठी सार्वत्रिक आहे. पण यामुळे
\olref[nou]{thm:no-univ} शी विरोध होतो.
\end{proof}

\begin{proof}
$h(e,x)$ संगणनक्षम आहे असे समजा. फलन $g$ ची पुढीलप्रमाणे व्याख्या करा:
\[
g(x) =
\begin{cases}
  0                & \text{जर $h(x,x) = 0$ असेल तर} \\
  \fundefined & \text{अन्यथा.}
\end{cases}
\]
फलन $g$ आंशिक संगणनक्षम आहे. उदाहरणार्थ, त्याची व्याख्या
$\umin{y}{h(x,x) = 0}$ अशी करता येते. म्हणून एखाद्या~$e$ साठी
प्रत्येक~$x$ बाबत $g(x) \simeq \fn{Un}(e,
x)$ असते. $g$ हे $e$ येथे परिभाषित आहे का? असल्यास, $g$ च्या
व्याख्येवरून $h(e,e) = 0$ ($h$ परिभाषित असेल तेव्हाच ते~$0$ मूल्य
घेऊ शकते). $h$ च्या व्याख्येनुसार, याचा अर्थ $\fn{Un}(e, e)$
अपरिभाषित आहे. पण प्रत्येक~$x$ साठी $g(x) \simeq \fn{Un}(e, x)$
या गृहितकावरून $g(e)$ अपरिभाषित ठरते—विरोध. दुसरीकडे, $g(e)$
अपरिभाषित असेल, तर $h(e,e) \neq 0$, म्हणून $h(e,e) = 1$.
त्यामुळे $\fn{Un}(e, e)$ परिभाषित आहे. पण
$g(x) \simeq \fn{Un}(e, x)$ असल्यामुळे $g(e)$ देखील परिभाषित
असेल. पुन्हा विरोध मिळतो.
\end{proof}

\begin{tagblock}{TMs}
\begin{explain}
हा युक्तिवाद ट्यूरिंग यंत्रांच्या भाषेत सांगता येतो. $H$ नावाचे
ट्यूरिंग यंत्र आहे असे समजा; ते ट्यूरिंग यंत्र~$E$ चे वर्णन आणि
आदान~$x$ घेते व $E$ हे आदान~$x$ वर थांबते का ते ठरवते.
मग आपण एकच आदान~$x$ घेणारे दुसरे ट्यूरिंग यंत्र~$G$ तयार करू शकतो.
ते निर्देशांक~$x$ असलेले यंत्र $M_x$ हे आदान~$x$ वर थांबते का हे
ठरवण्यासाठी $H$ चालवते आणि त्याच्या उलट वागते. म्हणजे, $M_x$
हे आदान~$x$ वर थांबते असे $H$ ने सांगितल्यास $G$ अनंत आवर्तनात
शिरते; आणि $M_x$ हे आदान~$x$ वर थांबत नाही असे $H$ ने
सांगितल्यास $G$ थांबते.
स्वतःचाच निर्देशांक आदान म्हणून दिल्यास $G$ थांबते का? वरील
युक्तिवाद दाखवतो की ते थांबते तेव्हाच थांबत नाही—हा विरोध आहे.
म्हणून असे ट्यूरिंग यंत्र~$H$ आहे हे गृहितक खोटे ठरते.
\end{explain}
\end{tagblock}

\end{document}
